\chapter{Formats in detail}\label{ch:formats}
\begin{goals}
\textbf{In this chapter} you learn what lamina does with each format beyond delimited text: Excel
workbooks, Datamine \file{.dm}, Apache Parquet and archives, including what each one can't do and why.
\end{goals}

Every format below becomes the same thing once it is open: a table of rows and named columns, each
column a number or text, which the grid, the filter, the calculated columns and the instruments treat
the same way. What differs is how the rows are reached, and so what costs what.

\section{Excel workbooks}\label{sec:xlsx}

Open an \file{.xlsx} or \file{.xlsm} and every sheet with data in it opens in a tab of its own, named
by the workbook and the sheet: \file{assays.xlsx › AU}. The first sheet is read at once; each of the
others is read the first time you look at its tab. Empty sheets get no tab (\cref{ch:views}).

The columns arrive typed as Excel typed them. A column is a number because the workbook stores numbers
in it, not because its text happened to parse, so a column of codes stored as text in Excel stays text,
leading zeros and all. The first row of the sheet is the header.

A workbook is a zip of XML that has to be decompressed and read whole, so unlike a CSV it is held in
memory, and lamina refuses one larger than 32~MB (the file's size on disk, compressed):
\code{workbook too large}. For a sheet that big, save it from Excel as CSV, and lamina windows the CSV at
any size.

\begin{tip}
Opening a sheet in lamina is often quicker than starting Excel, and the distribution over each header
says more about a column at a glance than scrolling it does.
\end{tip}

\section{Datamine .dm}\label{sec:dm}

A Datamine \file{.dm} file is a binary table: a page of field definitions, then fixed-length records.
lamina reads it directly, at any size, with no conversion: the header from the first page, then each
record by its position. Both Datamine precisions are read (single, with 2048-byte pages and 4-byte
numbers, and extended, with 4096-byte pages and 8-byte numbers), in either byte order, with the
extended format's long field names.

\begin{itemize}
\item \textbf{Numeric fields} are numbers; \textbf{alphanumeric fields} are text, however many words
they span.
\item \textbf{File constants}, fields Datamine stores once for the whole file (often \col{XINC},
\col{YINC}, \col{ZINC} in a regular model), appear as a column with the same value in every row.
\item \textbf{The missing-value sentinel} ($\pm10^{30}$, Datamine's ``absent'') reads as blank, not as
a huge number, so it doesn't enter means and sums (\cref{sec:blanks}).
\end{itemize}

Because every record is the same length and every field sits at the same offset within it, lamina can
read one column without decoding the others. A filter on a \file{.dm} reads only the fields its
expression names (\cref{sec:strided}); on a 17-column model filtering on one grade, that is most of the
file left unread.

\begin{caution}
A single-precision \file{.dm} stores each number in four bytes, which hold about seven significant
digits. lamina currently shows and exports those numbers with all the digits of their exact binary
value, so a grade stored as 0.1 appears as \code{0.10000000149011612}. The value is right to the
precision the file holds; the extra digits are noise. Set a number format to read them
(\cref{sec:formats}), and expect the noise in a CSV exported from such a file until this is fixed.
\end{caution}

\begin{caution}
A \file{.dm} inside an archive (\file{model.dm.gz}, or a zip holding a \file{.dm}) currently opens as a
binary file: lamina's archive path doesn't yet recognise Datamine. Decompress it first and open the
\file{.dm} itself.
\end{caution}

\section{Apache Parquet}\label{sec:parquet}

Parquet \parencite{parquetformat} is a columnar format: the file is cut into \term{row groups}, each
storing its columns one after another, compressed, with a footer at the end that describes the schema
and, for each row group and column, its minimum, maximum and count of missing values.

lamina reads the footer first. The schema comes from it, so a column is numeric because the file says
so. Then it reads only what it needs: the rows on screen, a few hundred at a time, or for a scan, only
the columns the scan uses.

\subsection{Types}

\begin{center}\small
\begin{tabularx}{\linewidth}{@{}lL@{}}
\toprule
\sffamily\bfseries in the file & \sffamily\bfseries in lamina \\
\midrule
integers and floating point & numbers \\
64-bit integers & numbers, when every value fits a JavaScript number exactly (up to $2^{53}$, about
$9\times10^{15}$); otherwise the whole column reads as text, so no digit is rounded away \\
text & text \\
booleans & the text \code{true} or \code{false} \\
dates and timestamps & ISO text: \code{2026-03-14}, or \code{2026-03-14T08:30:00Z} with a time \\
raw bytes & their hexadecimal text \\
nested columns (struct, list, map) & left out; the popover's facts name them under \ui{left out
(nested)} \\
\bottomrule
\end{tabularx}
\end{center}

\subsection{Compression}

lamina reads Parquet compressed with Snappy, gzip or Zstandard, or not compressed. It refuses LZ4 and
Brotli by name, before reading any data: \code{compressed with LZ4 — lamina reads Snappy, gzip, zstd or
uncompressed Parquet}. Their decoders would mean shipping WebAssembly, which lamina doesn't
(\cref{sec:no-wasm}). Most writers default to Snappy; in pandas, \code{df.to\_parquet(path)} writes
Snappy unless told otherwise.

\subsection{Row groups and pruning}\label{sec:pruning}

A filter on a Parquet file looks at each row group's footer statistics before reading it. If the
minimum and maximum of the columns the filter uses show that no row in the group can match, the group is
skipped unread. On a file sorted or spatially ordered by the filtered column, that can be almost the
whole file. The host API reports how many groups a filter skipped (\cref{app:host}); in the interface,
the effect is that the filter is fast.

The skip is conservative: a group is skipped only when the statistics prove no row matches. A filter
lamina can't reason about from minimum and maximum alone (a \code{matches} pattern, a calculated column)
reads every group.

\begin{tryit}
See what a Parquet file says about itself.
\begin{enumerate}
\item Write a Parquet file from pandas, or take one you have: \code{df.to\_parquet("model.parquet")}.
\item Open it in lamina and click the kind badge. The facts give the format, the column count, the
number of row groups and the compression.
\item Filter on a column the file is sorted by. Compare how quickly it answers with a filter on a column
it isn't sorted by.
\end{enumerate}
\end{tryit}

\section{Archives}\label{sec:archives}

lamina looks inside an archive and opens an entry without writing anything to disk. What it can do
depends on the compression, because some decompressors can only start at the beginning and go forward.

\begin{center}\small
\begin{tabularx}{\linewidth}{@{}lL@{}}
\toprule
\sffamily\bfseries archive & \sffamily\bfseries how it opens \\
\midrule
\file{.zip} & up to 32~MB: decompressed in memory, any entry. Larger: the entry list is read by streaming
through the file, and the chosen entry is read as a stream (below), at any size \\
\file{.gz} & up to 32~MB: decompressed in memory. Larger: read as a stream, at any size \\
\file{.tar} & up to 32~MB, in memory; a \file{.tar.gz}, \file{.tar.zst} and so on, up to 32~MB,
decompress first, then list their entries \\
\file{.zst}, \file{.bz2}, \file{.xz} & up to 32~MB, in memory; larger ones are refused (\code{archive
too large}), since the browser has no streaming decoder for them \\
\bottomrule
\end{tabularx}
\end{center}

An archive with one entry opens it; one with several lists them, with their sizes, and you pick. The
table is named by the archive and the entry: \file{survey.zip › collars.csv}.

\subsection{Reading a compressed stream}\label{sec:tape}

A large gzip file or zip entry can't be read from the middle: the decompressor has to start at the
beginning. lamina reads it through a \term{tape}: one decompressor reading forward, with the most recent
stretch of output kept in a buffer. Rows near where the tape is are served from the buffer; rows further
on are reached by reading forward; and rows behind it are reached by starting the decompressor again
from the top.

So scrolling forward through a 5~GB \file{.csv.gz} is cheap, a filter or a statistic (one forward pass)
costs what it would on the uncompressed file plus the decompression, and jumping back to the start of a
huge compressed file costs a re-read from the beginning. Its index is cached like any other file's
(\cref{sec:index-cache}).

\begin{why}
Decompressing a large archive to memory would bring back the limit windowing removes; decompressing it to
disk would need write access lamina doesn't ask for. The tape is the third way: nothing held whole, the
price paid only when you go backwards. If you will go back and forth through a huge compressed file
many times, decompress it once outside lamina and open the plain file.
\end{why}

Entries are recognised as their own formats once decompressed: a CSV inside a zip is a CSV, a Parquet
inside one is a Parquet (read from memory). A Datamine file inside an archive isn't yet (the caution
above).
